% ============================================================
% 第 4 章 场景与 ECS：EnTT 与 façade 架构
% 锚点：Engine/src/Scene/Scene.h, Components.h, Entity.h,
%       SceneHierarchyService.cpp, WorldTransformService.cpp,
%       Runtime/SceneRuntimeCoordinator, Systems/
% 时间线：8150251(2026-02-21) → 4dab189 → a6676de → 83dcf39
% ============================================================
\section{问题引入：对象之间说不清的关系}

有了渲染能力后，下一个问题是：场景里的东西怎么组织？一个``点光源挂在
移动的汽车上''的需求，就同时牵出了实体表示、组件组合、父子变换传播、
按类型遍历四件事。用面向对象的实体继承树（GameObject $\to$
LightObject/PhysicsObject……）很快会陷入菱形继承泥潭——这正是 ECS
（Entity-Component-System）要解决的问题。

本项目在首个提交日（2026-02-21）就选定了 EnTT 作为 ECS 底座
（提交 \texttt{8150251} ``Phase 2 Scene/ECS 场景系统 + 编辑器''），
此后经历三轮架构演进：运行时资源剥离（\texttt{4dab189}）、façade 收束
（\texttt{a6676de}）、组件元数据统一（\texttt{83dcf39}）。本章沿这条
演化线展开。

\section{通用原理：ECS 与 sparse-set 存储}

ECS 的核心三分法：\emph{实体}是 ID，\emph{组件}是纯数据 POD，
\emph{系统}是操作特定组件组合的逻辑。数据导向设计让缓存命中率成为
一等公民。

EnTT 的存储方案是 sparse-set：每个组件类型维护一个``稀疏数组（实体 ID
$\to$ 稠密下标）+ 稠密数组（紧凑的组件数据）''对。迭代某类组件时只走稠密
数组，内存连续、缓存友好；查询``拥有 A 且 B 的实体''则取两个集合的交集。
代价是非连续访问（如随机查询单个实体）需要一次稀疏跳转。理解这一点，
才能理解为什么本项目的渲染系统全部写成
\texttt{registry.view<TransformComponent, MeshRendererComponent>()}
式的批量遍历。

\section{本项目的设计与决策}

\subsection{组件即 POD}

\begin{lstlisting}
// Components.h —— 组件是纯数据 + 默认/拷贝构造
struct TransformComponent
{
    glm::vec3 Translation = {0.0f, 0.0f, 0.0f};
    glm::vec3 Rotation    = {0.0f, 0.0f, 0.0f};
    glm::vec3 Scale       = {1.0f, 1.0f, 1.0f};

    // 获取本地变换矩阵（TR S 复合）
    glm::mat4 GetTransform() const
    {
        glm::mat4 rotation = glm::toMat4(glm::quat(Rotation));
        return glm::translate(glm::mat4(1.0f), Translation)
             * rotation * glm::scale(glm::mat4(1.0f), Scale);
    }
};

struct RelationshipComponent
{
    UUID              ParentID = 0;   // 父实体 UUID，0 = 根节点
    std::vector<UUID> Children;       // 子实体 UUID 列表
};
\end{lstlisting}

两条纪律贯穿全部 20 余个组件：必须有默认构造和拷贝构造（序列化与 Undo
快照依赖值语义）；跨实体引用一律存 UUID 而非指针或 entt::entity
（RelationshipComponent 的 ParentID/Children 就是例证）。UUID 是唯一能
在``保存 $\to$ 重载 $\to$ 快照恢复''全流程中稳定存活的标识符。

\subsection{façade 化：Scene 是门面，不是仓库}

最初的 Scene 是个胖类：注册表、层级逻辑、世界变换计算、运行时资源、
音频视频容器全都塞在里面。三轮重构后它瘦成了六成员 façade
（Engine/src/Scene/Scene.h）：

\begin{lstlisting}
class Scene
{
public:
    Entity CreateEntity(const std::string& name = "Entity");
    void SetParent(Entity child, Entity parent);   // 委托 HierarchyService
    glm::mat4 GetWorldTransform(Entity entity);    // 委托 WorldTransformService
    void OnRuntimeStart();                         // 委托 RuntimeCoordinator
    entt::registry& GetRegistry() { return m_Registry; }
private:
    entt::registry m_Registry;
    SceneEntityIndex m_EntityIndex;          // UUID -> entity O(1) 索引
    WorldTransformCache m_TransformCache;
    ResourceLifecycleCoordinator m_LifecycleCoordinator;
    SceneEnvironmentState m_EnvironmentState;
    std::unique_ptr<SceneRuntimeCoordinator> m_RuntimeCoordinator;
};
\end{lstlisting}

拆出去的逻辑去哪了？变成无状态静态服务类。\texttt{WorldTransformService}
沿 RelationshipComponent 递归上溯父链复合世界矩阵：

\begin{lstlisting}
// WorldTransformService.cpp —— 递归 + 缓存 + 深度防御
glm::mat4 ComputeWorldTransformImpl(reg, entity, index,
                                    cache, depth)
{
    if (depth >= kMaxDepth)   // kMaxDepth=64
    {
        ENGINE_CORE_WARN("... possible cycle in hierarchy");
        return reg.get<TransformComponent>(entity).GetTransform();
    }
    if (cache && cache->TryGet(entity, cached))
        return cached;

    auto& result = reg.get<TransformComponent>(entity).GetTransform();
    if (reg.all_of<RelationshipComponent>(entity))
    {
        auto& rel = reg.get<RelationshipComponent>(entity);
        if (rel.ParentID != 0) {
            entt::entity parent = index.Find(rel.ParentID);
            if (parent != entt::null && reg.valid(parent))
                result = ComputeWorldTransformImpl(
                    reg, parent, index, cache, depth + 1) * result;
        }
    }
    if (cache) cache->Put(entity, result);
    return result;
}
\end{lstlisting}

三个防御细节值得注意：\textbf{kMaxDepth=64 兜底循环层级}（数据损坏时不至
栈溢出）；\textbf{缓存命中短路递归}（编辑器每帧对全场景求世界变换）；而
\texttt{SetParent} 时会先记下子实体的旧世界矩阵，换父后用
$\mathrm{local} = P_{new}^{-1} \cdot W_{old}$ 反解本地变换——保证``拖拽换父
时物体在世界空间里纹丝不动''（Engine/src/Scene/SceneHierarchyService.cpp，
配合 GLM \texttt{matrix\_decompose} 分解回 T/R/S）。

\subsection{双场景快照：Play 按钮的安全网}

点击 Play 会进入物理与脚本的运行态，但编辑器必须能在 Stop 后回到进入前的
精确状态。实现方式不是记录 diff，而是\textbf{整场复制}：EditorSceneSession
维护编辑态与运行态两个 Scene 实例（重构提交 \texttt{b8bceb1}），Play 时把
编辑场景 Copy 一份供模拟使用，Stop 时直接丢弃运行副本。OpenAL Source、
FFmpeg 解码器这类无法安全复制的原生资源由 RuntimeStore 在 Stop 时强制回收
（Phase 1 提交 \texttt{4dab189} 的``运行时资源剥离''正是为此铺路）。
粗暴、简单、可靠——快照成本对一个编辑器会话完全可以接受。

\subsection{Systems 目录：渲染逻辑的静态归拢}

MeshRenderSystem / ShadowSystem / SkyboxSystem / LightSystem /
TerrainRenderSystem / GrassRenderSystem / AudioSystem / VideoSystem 八个
静态系统类构成 Systems/ 目录。它们没有状态、不互相调用，入口全是静态方法，
通过 registry view 批量处理组件。``系统无状态''让渲染管线的执行顺序完全
由 SceneRenderer 编排层决定（第 3 章 8-Pass），职责切分干净。

\section{实现走读:删除一个实体有多难}

DestroyEntity 要处理：递归删除整个子树、清理 UUID 索引、触发资源回调、
从父节点的 Children 列表移除自己。最初实现里``找子树''用的是对每个节点
做全场景扫描，复杂度 $O(N^2)$；优化版（\texttt{641f8ae}）先一次性收集
子树再批量删除。另一个隐蔽 bug 是 UUID 索引允许重复插入
（\texttt{595e461} 修复为检测并拒绝），否则 Find 返回歧义结果，拾取与
序列化都会错乱。EnTT 自身也有兼容性暗礁——\texttt{size\_hint} API 在
版本间行为变化导致实体计数错误（\texttt{2490c50} 改用兼容写法）。

\section{踩坑记录}

\begin{description}
  \item[组件里藏指针] 早期曾有组件直接持有裸指针，Play/Stop 快照恢复后
        变悬垂。教训固化为纪律：组件只存 POD 与 UUID，重资源走 AssetHandle
        （第 6 章）。
  \item[层级环] 手滑把自己设为自己的祖先，世界变换递归立即爆栈。
        kMaxDepth 防御 + IsAncestorOf 前置检查双重保险。
  \item[O(N²) 删除] 大场景删除根节点卡顿数秒，profiler 定位到逐节点
        全表扫描；先收集后删降到 $O(N)$（641f8ae）。
\end{description}

\section{小结与延伸思考}

本章的演化路径——胖 Scene $\to$ façade + 无状态服务 + 协调器——本质是
把``数据（registry）/ 规则（service）/ 流程（coordinator）''三种变化
频率不同的东西分开放。它带来的副产品是可测试性：HierarchyService、
WorldTransformService、EntityIndex 都因此获得了独立单元测试
（tests/ 下 TestSceneHierarchy/TestWorldTransform/TestSceneEntityIndex）。

延伸思考：EnTT 的 sparse-set 对``按类型批量遍历''最优，但对``按实体随机
访问多组件''并不友好。如果未来要做大规模开放世界流式加载，可能需要在
其上再建一层空间索引。当前项目规模下，简单性红利仍然大于替代方案的收益。
